\section{Data Processing}
\label{sec:supp-data}

\subsection{Data selection}
\label{subsec:supp-data-sources}
Our selection of datasets for building \LaViDa\ is detailed in Tab.~\ref{tab:lavida-details}. 
This collection is intended to provide images covering well various downstream vision tasks both for image-level and dense recognition.

\subsection{Image similarity}
\label{subsec:supp-data-similarity}
We employ cosine similarity to compare image features (whether ours or feature generated for deduplication) with the following similarity function $m$:
$$m(s, r) = \text{cosine-similarity}\left(f\left(s\right),f\left(r\right)\right) = \frac{f(s)\cdot{}f(r)}{\lVert f(s)\rVert_2\lVert f(r)\rVert_2}$$
where $s$ and $r$ are a pair of images to compare and $f$ is the model generating features.

\subsection{Deduplication}
\label{subsec:supp-data-dedup}

\paragraph{Self-deduplication.} 
To deduplicate our uncurated data source of 1.3B images, we compute and use the embeddings generated by \cite{pizzi2022self} and retrieve the $k = 64$ nearest neighbors of each image (using cosine similarity). 
Considering only neighbors with a similarity ${>}0.6$, we extract the connected components of the associated $k$-NN graph thanks to a scalable disjoint set data structure implementation. 
We then only keep one representative for each component of duplicate images. 
This results in a self-deduplicated data source of 1.1B images.

\paragraph{Relative deduplication}
To reduce redundancy and also properly evaluate the performance of our features, we discard remaining images of our self-deduplicated data source that are too similar to train and test splits of our evaluation datasets. 
To achieve this, we apply a similar procedure as for self-deduplication, with a stricter similarity ${>}0.45$, this time identifying the duplicate components (if any) to which each reference image belong and discarding it entirely. 
This results in a self- and relatively-deduplicated data source of 744M images.

\subsection{Retrieval}
\label{subsec:supp-data-retrieval}
We employ two approaches to augment dataset via retrieval: sample-based and cluster-based. 
The first one, sample-based, applies to datasets larger than 1M images and consists in collecting a fixed number $k$ of nearest images for each sample image of the dataset to retrieve, effectively trying to multiply by $k$ the size of the dataset. 
We use $k = 4$ for Google Landmarks v2 and ImageNet-22k but a larger $k = 32$ to make this specific retrieval a core part of our \LaViDa{} dataset. 
For smaller datasets, the second approach, cluster-based, consists in first clustering our uncurated data source into $100,000$ separate clusters thanks to a distributed $k$-means implementation. 
Each cluster should capture different types of image concept and contents. 
We then pick $10,000$ images from each cluster associated with more than $3$ images of the retrieved dataset. 
As this can result in a very large number of retrieved images for some dataset, we restrict such retrievals to a maximum of 1M images to maintain the balance between the different datasets within \LaViDa.

\begin{table*}
    \centering
    \resizebox{\linewidth}{!}{
    \begin{tabular}{@{} l l r c r r @{}}
    \toprule
    Task                    & Dataset / Split                     & Images     & Retrieval & Retrieved  & Final  \\
    \midrule
    classification          & ImageNet-22k / --                   & 14,197,086 & as is     & --         & 14,197,086 \\ 
    classification          & ImageNet-22k / --                   & 14,197,086 & sample    & 56,788,344 & 56,788,344 \\ 
    classification          & ImageNet-1k / train                 &  1,281,167 & sample    & 40,997,344 & 40,997,344 \\ 
    \midrule
    fine-grained classif.   & Caltech 101 / train                 &      3,030 & cluster   &  2,630,000 &  1,000,000 \\ 
    fine-grained classif.   & CUB-200-2011 / train                &      5,994 & cluster   &  1,300,000 &  1,000,000 \\ 
    fine-grained classif.   & DTD / train1                        &      1,880 & cluster   &  1,580,000 &  1,000,000 \\ 
    fine-grained classif.   & FGVC-Aircraft / train               &      3,334 & cluster   &  1,170,000 &  1,000,000 \\ 
    fine-grained classif.   & Flowers-102 / train                 &      1,020 & cluster   &  1,060,000 &  1,000,000 \\ 
    fine-grained classif.   & Food-101 / train                    &     75,750 & cluster   & 21,670,000 &  1,000,000 \\ 
    fine-grained classif.   & Oxford-IIIT Pet / trainval          &      3,680 & cluster   &  2,750,000 &  1,000,000 \\ 
    fine-grained classif.   & Stanford Cars / train               &      8,144 & cluster   &  7,220,000 &  1,000,000 \\ 
    fine-grained classif.   & SUN397 / train1                     &     19,850 & cluster   & 18,950,000 &  1,000,000 \\ 
    fine-grained classif.   & Pascal VOC 2007 / train             &      2,501 & cluster   &  1,010,000 &  1,000,000 \\ 
    \midrule
    segmentation            & ADE20K / train                      &     20,210 & cluster   & 20,720,000 &  1,000,000 \\ 
    segmentation            & Cityscapes / train                  &      2,975 & cluster   &  1,390,000 &  1,000,000 \\ 
    segmentation            & Pascal VOC 2012 (seg.) / trainaug   &      1,464 & cluster   & 10,140,000 &  1,000,000 \\ 
    \midrule
    depth estimation        & Mapillary SLS / train               &  1,434,262 & as is     & --         &  1,434,262 \\ 
    depth estimation        & KITTI / train (Eigen)               &     23,158 & cluster   &  3,700,000 &  1,000,000 \\ 
    depth estimation        & NYU Depth V2 / train                &     24,231 & cluster   & 10,850,000 &  1,000,000 \\ 
    depth estimation        & SUN RGB-D / train                   &      4,829 & cluster   &  4,870,000 &  1,000,000 \\ 
    \midrule
    retrieval               & Google Landmarks v2 / train (clean) &  1,580,470 & as is     & --         &  1,580,470 \\ 
    retrieval               & Google Landmarks v2 / train (clean) &  1,580,470 & sample    &  6,321,880 &  6,321,880 \\ 
    retrieval               & AmsterTime / new                    &      1,231 & cluster   &    960,000 &    960,000 \\ 
    retrieval               & AmsterTime / old                    &      1,231 & cluster   &    830,000 &    830,000 \\ 
    retrieval               & Met / train                         &    397,121 & cluster   & 62,860,000 &  1,000,000 \\ 
    retrieval               & Revisiting Oxford / base            &      4,993 & cluster   &  3,680,000 &  1,000,000 \\ 
    retrieval               & Revisiting Paris / base             &      6,322 & cluster   &  3,660,000 &  1,000,000 \\ 
    \bottomrule
                            &                                     &            &           &            & 142,109,386 \\
    \end{tabular}
    }
    \caption{
      \textbf{Composition of our \LaViDa{} dataset.} 
      We report the list of datasets and associated splits used to build the dataset, how they were included (as is without retrieval or via sample-based or cluster-based retrieval). 
      For retrievals, we indicate the actual number of retrieved images and the final number included in the dataset. 
      \rthree{
          We chose to include as many datasets as possible in the pretraining data in order to cover as many domains as possible. 
          We kept a few datasets aside in order to evaluate performance outside of the pretraining domain. 
          More details about dataset usages can be found in Table~\ref{tab:datasets_list}.
      }
      }
    \label{tab:lavida-details}
\end{table*}

